\documentclass[10pt,twocolumn,letterpaper]{article}
\usepackage[review]{cvpr}
%% This file contains a number of tweaks that are typically applied to the main document.
%% They are not enabled by default, but can be enabled by uncommenting the relevant lines.


%%
%% Official support for 2-column teaser figure
%%
\usepackage{cuted} %< for strip environment

%%
%% For auto-referencing of figures (see fig/teaser.tex and \label{} therein) 
%%
\usepackage{currfile} %< for \currfilebase macro

%%
%% caption and (global) spacing overrides
%%
\usepackage{caption} %< for "\captionof" commands (teaser)
% \setlength{\abovecaptionskip}{.5em} % space above caption
% \setlength{\belowcaptionskip}{1em} % space below caption

%%
%% Inline annotations; for predefined colors, refer to "dvipsnames" in the xcolor package:
%% https://tinyurl.com/overleaf-colors
%%
\newcommand{\red}[1]{{\color{red}#1}}
\newcommand{\todo}[1]{{\color{red}#1}}
\newcommand{\TODO}[1]{\textbf{\color{red}[TODO: #1]}}
%%
%% Disable for camera ready / submission by uncommenting these lines:
%%
% \renewcommand{\TODO}[1]{}
% \renewcommand{\todo}[1]{#1}

%%
%% Work harder in optimizing text layout. Typically shrinks text by 1/6 of page, enable
%% it at the very end of the writing process, when you are just above the page limit.
%%
% \usepackage{microtype}

%%
%% Fine-tune paragraph spacing.
%%
% \renewcommand{\paragraph}[1]{\vspace{.5em}\noindent\textbf{#1.}}

%%
%% Globally adjusts space between figure and caption.
%%
% \setlength{\abovecaptionskip}{.5em}

%%
%% Allows "the use of \paper to refer to the project name"
%% with automatic management of space at the end of the word.
%%
% \usepackage{xspace}
% \newcommand{\paper}{ProjectName\xspace}

%%
%% Commonly used math definitions.
%%
% \DeclareMathOperator*{\argmin}{arg\,min}
% \DeclareMathOperator*{\argmax}{arg\,max}

%%
%% Tighten underline.
%%
% \usepackage{soul}
% \setuldepth{foobar}

%%
%% Reconfigure paragraph (tweaks spacing)
%%

\usepackage{iftex}
\ifXeTeX\usepackage{fontspec}
\setmainfont{texgyretermes}[Extension=.otf,UprightFont=*-regular,BoldFont=*-bold,ItalicFont=*-italic,BoldItalicFont=*-bolditalic]
\setsansfont{texgyreheros}[Extension=.otf,UprightFont=*-regular,BoldFont=*-bold,ItalicFont=*-italic,BoldItalicFont=*-bolditalic]
\setmonofont{texgyrecursor}[Extension=.otf,UprightFont=*-regular,BoldFont=*-bold,ItalicFont=*-italic,BoldItalicFont=*-bolditalic]
\fi
\usepackage{booktabs,amsmath,amssymb,graphicx,xspace,makecell}
\definecolor{cvprblue}{rgb}{0.21,0.49,0.74}
\usepackage[breaklinks,colorlinks,allcolors=cvprblue]{hyperref}
\raggedbottom

\def\paperID{*****}
\def\confName{ICCV}
\def\confYear{2027}
\newcommand{\ours}{S2M\xspace}
\renewcommand{\thesection}{S\arabic{section}}
\renewcommand{\thetable}{S\arabic{table}}
% generated by scripts/paper_assets.py from the TEST records -- do not edit
\newcommand{\nEBThree}{55.0}
\newcommand{\nEBThreeLo}{49.3}
\newcommand{\nEBThreeHi}{60.7}
\newcommand{\nEBThreeBase}{66.7}
\newcommand{\nEBThreeCommon}{65.0}
\newcommand{\nEBThreeComplex}{65.0}
\newcommand{\nEBThreeVisual}{53.3}
\newcommand{\nEBThreeLong}{25.0}
\newcommand{\nEBThreeAWQ}{48.0}
\newcommand{\nEBThreeAWQLo}{42.3}
\newcommand{\nEBThreeAWQHi}{53.7}
\newcommand{\nEBThreeAWQBase}{65.0}
\newcommand{\nEBThreeAWQCommon}{65.0}
\newcommand{\nEBThreeAWQComplex}{60.0}
\newcommand{\nEBThreeAWQVisual}{38.3}
\newcommand{\nEBThreeAWQLong}{11.7}
\newcommand{\nSThree}{79.4}
\newcommand{\nSThreeLo}{75.0}
\newcommand{\nSThreeHi}{83.7}
\newcommand{\nSThreeBase}{89.6}
\newcommand{\nSThreeCommon}{84.2}
\newcommand{\nSThreeComplex}{88.3}
\newcommand{\nSThreeVisual}{52.9}
\newcommand{\nSThreeLong}{82.1}
\newcommand{\nRGBThree}{76.5}
\newcommand{\nRGBThreeLo}{71.8}
\newcommand{\nRGBThreeHi}{81.0}
\newcommand{\nRGBThreeBase}{85.0}
\newcommand{\nRGBThreeCommon}{80.8}
\newcommand{\nRGBThreeComplex}{85.8}
\newcommand{\nRGBThreeVisual}{52.5}
\newcommand{\nRGBThreeLong}{78.3}
\newcommand{\nEBTwo}{20.3}
\newcommand{\nEBTwoLo}{16.0}
\newcommand{\nEBTwoHi}{25.0}
\newcommand{\nEBTwoBase}{21.7}
\newcommand{\nEBTwoCommon}{21.7}
\newcommand{\nEBTwoComplex}{28.3}
\newcommand{\nEBTwoVisual}{20.0}
\newcommand{\nEBTwoLong}{10.0}
\newcommand{\nEBTwoAWQ}{6.3}
\newcommand{\nEBTwoAWQLo}{3.7}
\newcommand{\nEBTwoAWQHi}{9.3}
\newcommand{\nEBTwoAWQBase}{6.7}
\newcommand{\nEBTwoAWQCommon}{6.7}
\newcommand{\nEBTwoAWQComplex}{6.7}
\newcommand{\nEBTwoAWQVisual}{8.3}
\newcommand{\nEBTwoAWQLong}{3.3}
\newcommand{\nEBTwoAWQraw}{5.7}
\newcommand{\nEBTwoAWQrawLo}{3.3}
\newcommand{\nEBTwoAWQrawHi}{8.3}
\newcommand{\nEBTwoAWQrawBase}{8.3}
\newcommand{\nEBTwoAWQrawCommon}{6.7}
\newcommand{\nEBTwoAWQrawComplex}{6.7}
\newcommand{\nEBTwoAWQrawVisual}{6.7}
\newcommand{\nEBTwoAWQrawLong}{0.0}
\newcommand{\nSTwo}{70.0}
\newcommand{\nSTwoLo}{64.7}
\newcommand{\nSTwoHi}{75.0}
\newcommand{\nSTwoBase}{78.3}
\newcommand{\nSTwoCommon}{71.7}
\newcommand{\nSTwoComplex}{71.7}
\newcommand{\nSTwoVisual}{53.3}
\newcommand{\nSTwoLong}{75.0}
\newcommand{\nOracle}{91.3}
\newcommand{\nOracleLo}{88.0}
\newcommand{\nOracleHi}{94.3}
\newcommand{\nOracleBase}{93.3}
\newcommand{\nOracleCommon}{93.3}
\newcommand{\nOracleComplex}{93.3}
\newcommand{\nOracleVisual}{93.3}
\newcommand{\nOracleLong}{83.3}
\newcommand{\nGainThree}{24.4}
\newcommand{\nGainThreeLo}{18.0}
\newcommand{\nGainThreeHi}{30.8}
\newcommand{\nGainThreeAbs}{24.4}
\newcommand{\nGainThreeP}{{<}10^{-5}}
\newcommand{\nGainThreeBase}{+22.9}
\newcommand{\nGainThreeCommon}{+19.2}
\newcommand{\nGainThreeComplex}{+23.3}
\newcommand{\nGainThreeVisual}{-0.4}
\newcommand{\nGainThreeLong}{+57.1}
\newcommand{\nGainTwo}{49.7}
\newcommand{\nGainTwoLo}{43.0}
\newcommand{\nGainTwoHi}{56.3}
\newcommand{\nGainTwoAbs}{49.7}
\newcommand{\nGainTwoP}{{<}10^{-5}}
\newcommand{\nGainTwoBase}{+56.7}
\newcommand{\nGainTwoCommon}{+50.0}
\newcommand{\nGainTwoComplex}{+43.3}
\newcommand{\nGainTwoVisual}{+33.3}
\newcommand{\nGainTwoLong}{+65.0}
\newcommand{\nGainRGB}{21.5}
\newcommand{\nGainRGBLo}{14.7}
\newcommand{\nGainRGBHi}{28.2}
\newcommand{\nGainRGBAbs}{21.5}
\newcommand{\nGainRGBP}{{<}10^{-5}}
\newcommand{\nGainRGBBase}{+18.3}
\newcommand{\nGainRGBCommon}{+15.8}
\newcommand{\nGainRGBComplex}{+20.8}
\newcommand{\nGainRGBVisual}{-0.8}
\newcommand{\nGainRGBLong}{+53.3}
\newcommand{\nGainThreeAWQ}{31.4}
\newcommand{\nGainThreeAWQLo}{24.8}
\newcommand{\nGainThreeAWQHi}{37.9}
\newcommand{\nGainThreeAWQAbs}{31.4}
\newcommand{\nGainThreeAWQP}{{<}10^{-5}}
\newcommand{\nGainThreeAWQBase}{+24.6}
\newcommand{\nGainThreeAWQCommon}{+19.2}
\newcommand{\nGainThreeAWQComplex}{+28.3}
\newcommand{\nGainThreeAWQVisual}{+14.6}
\newcommand{\nGainThreeAWQLong}{+70.4}
\newcommand{\nGainTwoAWQ}{63.7}
\newcommand{\nGainTwoAWQLo}{58.0}
\newcommand{\nGainTwoAWQHi}{69.3}
\newcommand{\nGainTwoAWQAbs}{63.7}
\newcommand{\nGainTwoAWQP}{{<}10^{-5}}
\newcommand{\nGainTwoAWQBase}{+71.7}
\newcommand{\nGainTwoAWQCommon}{+65.0}
\newcommand{\nGainTwoAWQComplex}{+65.0}
\newcommand{\nGainTwoAWQVisual}{+45.0}
\newcommand{\nGainTwoAWQLong}{+71.7}
\newcommand{\nQuantThree}{7.0}
\newcommand{\nQuantThreeLo}{0.7}
\newcommand{\nQuantThreeHi}{13.3}
\newcommand{\nQuantThreeAbs}{7.0}
\newcommand{\nQuantThreeP}{{=}0.04}
\newcommand{\nQuantThreeBase}{+1.7}
\newcommand{\nQuantThreeCommon}{+0.0}
\newcommand{\nQuantThreeComplex}{+5.0}
\newcommand{\nQuantThreeVisual}{+15.0}
\newcommand{\nQuantThreeLong}{+13.3}
\newcommand{\nQuantTwo}{14.0}
\newcommand{\nQuantTwoLo}{8.7}
\newcommand{\nQuantTwoHi}{19.3}
\newcommand{\nQuantTwoAbs}{14.0}
\newcommand{\nQuantTwoP}{{<}10^{-5}}
\newcommand{\nQuantTwoBase}{+15.0}
\newcommand{\nQuantTwoCommon}{+15.0}
\newcommand{\nQuantTwoComplex}{+21.7}
\newcommand{\nQuantTwoVisual}{+11.7}
\newcommand{\nQuantTwoLong}{+6.7}
\newcommand{\nGapRGB}{-2.9}
\newcommand{\nGapRGBLo}{-6.3}
\newcommand{\nGapRGBHi}{0.3}
\newcommand{\nGapRGBAbs}{2.9}
\newcommand{\nGapRGBP}{{=}0.10}
\newcommand{\nGapRGBBase}{-4.6}
\newcommand{\nGapRGBCommon}{-3.3}
\newcommand{\nGapRGBComplex}{-2.5}
\newcommand{\nGapRGBVisual}{-0.4}
\newcommand{\nGapRGBLong}{-3.8}
\newcommand{\nGapOracle}{11.9}
\newcommand{\nGapOracleLo}{7.3}
\newcommand{\nGapOracleHi}{16.5}
\newcommand{\nGapOracleAbs}{11.9}
\newcommand{\nGapOracleP}{{<}10^{-5}}
\newcommand{\nGapOracleBase}{+3.8}
\newcommand{\nGapOracleCommon}{+9.2}
\newcommand{\nGapOracleComplex}{+5.0}
\newcommand{\nGapOracleVisual}{+40.4}
\newcommand{\nGapOracleLong}{+1.2}
\newcommand{\nMcNThreeWorst}{{=}8.7{\times}10^{-12}}
\newcommand{\nMcNTwoWorst}{{=}2.3{\times}10^{-33}}
\newcommand{\nMcNRGBWorst}{{=}5.2{\times}10^{-9}}
\newcommand{\nDiscTwoWins}{163}
\newcommand{\nDiscTwoLosses}{14}
\newcommand{\nDiscTwoP}{{=}2.3{\times}10^{-33}}
\newcommand{\nDiscThreeWins}{95}
\newcommand{\nDiscThreeLosses}{21}
\newcommand{\nDiscThreeP}{{=}2.0{\times}10^{-12}}
\newcommand{\nDiscGlobalWins}{59}
\newcommand{\nDiscGlobalLosses}{65}
\newcommand{\nDiscGlobalP}{{=}0.65}
\newcommand{\nRGBGlobal}{53.0}
\newcommand{\nGainGlobal}{-2.0}
\newcommand{\nVisDiffLo}{-15.0}
\newcommand{\nVisDiffHi}{+14.6}
\newcommand{\nVisDiffP}{{=}1.00}
\newcommand{\nQuantThreeNamedMax}{+5.0}
\newcommand{\nQuantThreeWins}{58}
\newcommand{\nQuantThreeLosses}{37}
\newcommand{\nQuantTwoWins}{58}
\newcommand{\nQuantTwoLosses}{16}
\newcommand{\nSThreeSD}{0.5}
\newcommand{\nRGBThreeSD}{0.2}
\newcommand{\nSThreeMin}{79.0}
\newcommand{\nSThreeMax}{80.0}
\newcommand{\nScaleDrop}{9.4}
\newcommand{\nScaleDropEB}{34.7}
\newcommand{\nOracleGapLong}{1.3}
\newcommand{\nStepN}{140}
\newcommand{\nStepS}{11.1}
\newcommand{\nStepEB}{12.6}
\newcommand{\nStepRed}{12}
\newcommand{\nStepFewer}{79}
\newcommand{\nStepMore}{38}
\newcommand{\nStepP}{{=}1.9{\times}10^{-4}}
\newcommand{\nTokS}{14.4}
\newcommand{\nTokSc}{513}
\newcommand{\nCallsS}{52.6}
\newcommand{\nTokEB}{27.1}
\newcommand{\nTokEBc}{2.3}
\newcommand{\nCallsEB}{7.8}
\newcommand{\nEBThreeParse}{1 of its 2,342 replies cannot be parsed and is replaced by a random action, as in the benchmark}
\newcommand{\nFailEBThree}{135}
\newcommand{\nNeverEBThree}{12}
\newcommand{\nLostEBThree}{43}
\newcommand{\nNearEBThree}{34}
\newcommand{\nBlockedEBThree}{26}
\newcommand{\nFailSThree}{61.8}
\newcommand{\nNeverSThree}{0.8}
\newcommand{\nLostSThree}{28.5}
\newcommand{\nNearSThree}{9.8}
\newcommand{\nBlockedSThree}{0.5}
\newcommand{\nNeverLongEBThree}{12}
\newcommand{\nRecallOtherLo}{93}
\newcommand{\nRecallOtherHi}{98}
\newcommand{\nRecallVisLo}{63}
\newcommand{\nRecallVisHi}{67}
\newcommand{\nVisFail}{113}
\newcommand{\nVisFailNoCand}{65}
\newcommand{\nVisSRLo}{50.0}
\newcommand{\nVisSRHi}{56.7}
\newcommand{\nProbeThreeBase}{78.7}
\newcommand{\nProbeThreeCommon}{69.3}
\newcommand{\nProbeThreeComplex}{77.3}
\newcommand{\nProbeThreeVisual}{49.3}
\newcommand{\nProbeTwoBase}{40.0}
\newcommand{\nProbeTwoCommon}{33.3}
\newcommand{\nProbeTwoComplex}{37.3}
\newcommand{\nProbeTwoVisual}{18.7}
\newcommand{\nDRmean}{0.10}
\newcommand{\nDRfinal}{0.15}
\newcommand{\nRangeSensor}{0.97}
\newcommand{\nRangeGlobal}{0.90}
\newcommand{\nRangeFloor}{0.93}
\newcommand{\nShortGlobal}{66}
\newcommand{\nShortFloor}{10.0}
\newcommand{\nProbeViews}{240}
\newcommand{\nBlockGlobal}{42.5}
\newcommand{\nBlockFloor}{9.9}
\newcommand{\nAbsRelGlobal}{0.134}
\newcommand{\nAbsRelFloor}{0.115}
\newcommand{\nShortSensorMax}{0}
\newcommand{\nVisGapShare}{68}
\newcommand{\nCapAWQ}{14.8}
\newcommand{\nIdFirstAWQ}{92.6}
\newcommand{\nIdLaterAWQ}{79.3}
\newcommand{\nCapFP}{3.2}
\newcommand{\nIdFirstFP}{99.6}
\newcommand{\nIdLaterFP}{95.7}
\newcommand{\nRoboN}{1800}
\newcommand{\nRoboSR}{59.6}
\newcommand{\nRoboSRLo}{57.3}
\newcommand{\nRoboSRHi}{61.8}
\newcommand{\nRoboSPL}{35.9}
\newcommand{\nRoboSPLLo}{34.2}
\newcommand{\nRoboSPLHi}{37.6}
\newcommand{\nRoboBest}{88.0}
\newcommand{\nRoboWorst}{25.3}
\newcommand{\nRoboErr}{17}
\newcommand{\nAblFusion}{-10.4}
\newcommand{\nAblFusionAbs}{10.4}
\newcommand{\nAblFusionLo}{-15.0}
\newcommand{\nAblFusionHi}{-6.0}
\newcommand{\nAblFusionP}{{=}8.1{\times}10^{-5}}
\newcommand{\nAblFusionBase}{-11.3}
\newcommand{\nAblFusionCommon}{-14.2}
\newcommand{\nAblFusionComplex}{-13.3}
\newcommand{\nAblFusionVisual}{-11.2}
\newcommand{\nAblFusionLong}{-2.1}
\newcommand{\nAblFusionOtherMin}{-14.2}
\newcommand{\nAblFusionOtherMax}{-11.2}
\newcommand{\nAblScan}{0.2}
\newcommand{\nAblScanAbs}{0.2}
\newcommand{\nAblScanLo}{-4.2}
\newcommand{\nAblScanHi}{4.8}
\newcommand{\nAblScanP}{{=}1.00}
\newcommand{\nAblScanBase}{+5.4}
\newcommand{\nAblScanCommon}{+7.5}
\newcommand{\nAblScanComplex}{+3.3}
\newcommand{\nAblScanVisual}{+18.8}
\newcommand{\nAblScanLong}{-33.7}
\newcommand{\nAblScanOtherMin}{+3.3}
\newcommand{\nAblScanOtherMax}{+18.8}
\newcommand{\nAblExhaust}{-1.1}
\newcommand{\nAblExhaustAbs}{1.1}
\newcommand{\nAblExhaustLo}{-2.5}
\newcommand{\nAblExhaustHi}{0.2}
\newcommand{\nAblExhaustP}{{=}0.89}
\newcommand{\nAblExhaustBase}{-1.3}
\newcommand{\nAblExhaustCommon}{+0.8}
\newcommand{\nAblExhaustComplex}{+0.0}
\newcommand{\nAblExhaustVisual}{-6.2}
\newcommand{\nAblExhaustLong}{+1.3}
\newcommand{\nAblExhaustOtherMin}{-6.2}
\newcommand{\nAblExhaustOtherMax}{+0.8}
\newcommand{\nAblNegative}{1.6}
\newcommand{\nAblNegativeAbs}{1.6}
\newcommand{\nAblNegativeLo}{-1.2}
\newcommand{\nAblNegativeHi}{4.5}
\newcommand{\nAblNegativeP}{{=}1.00}
\newcommand{\nAblNegativeBase}{+3.7}
\newcommand{\nAblNegativeCommon}{-0.8}
\newcommand{\nAblNegativeComplex}{+0.0}
\newcommand{\nAblNegativeVisual}{+5.4}
\newcommand{\nAblNegativeLong}{-0.4}
\newcommand{\nAblNegativeOtherMin}{-0.8}
\newcommand{\nAblNegativeOtherMax}{+5.4}
\newcommand{\nAblTiles}{-1.1}
\newcommand{\nAblTilesAbs}{1.1}
\newcommand{\nAblTilesLo}{-4.5}
\newcommand{\nAblTilesHi}{2.4}
\newcommand{\nAblTilesP}{{=}1.00}
\newcommand{\nAblTilesBase}{-2.9}
\newcommand{\nAblTilesCommon}{-4.2}
\newcommand{\nAblTilesComplex}{-3.3}
\newcommand{\nAblTilesVisual}{+7.1}
\newcommand{\nAblTilesLong}{-2.1}
\newcommand{\nAblTilesOtherMin}{-4.2}
\newcommand{\nAblTilesOtherMax}{+7.1}
\newcommand{\nAblZoom}{-2.8}
\newcommand{\nAblZoomAbs}{2.8}
\newcommand{\nAblZoomLo}{-5.7}
\newcommand{\nAblZoomHi}{0.1}
\newcommand{\nAblZoomP}{{=}0.57}
\newcommand{\nAblZoomBase}{-4.6}
\newcommand{\nAblZoomCommon}{-4.2}
\newcommand{\nAblZoomComplex}{-5.0}
\newcommand{\nAblZoomVisual}{-1.2}
\newcommand{\nAblZoomLong}{+1.3}
\newcommand{\nAblZoomOtherMin}{-5.0}
\newcommand{\nAblZoomOtherMax}{-1.2}
\newcommand{\nAblStrongest}{1.2}
\newcommand{\nAblStrongestAbs}{1.2}
\newcommand{\nAblStrongestLo}{-0.3}
\newcommand{\nAblStrongestHi}{2.9}
\newcommand{\nAblStrongestP}{{=}0.89}
\newcommand{\nAblStrongestBase}{+0.4}
\newcommand{\nAblStrongestCommon}{+0.8}
\newcommand{\nAblStrongestComplex}{+1.7}
\newcommand{\nAblStrongestVisual}{+2.1}
\newcommand{\nAblStrongestLong}{+1.3}
\newcommand{\nAblStrongestOtherMin}{+0.4}
\newcommand{\nAblStrongestOtherMax}{+2.1}
\newcommand{\nAblOdom}{-0.1}
\newcommand{\nAblOdomAbs}{0.1}
\newcommand{\nAblOdomLo}{-3.1}
\newcommand{\nAblOdomHi}{2.9}
\newcommand{\nAblOdomP}{{=}1.00}
\newcommand{\nAblOdomBase}{+2.1}
\newcommand{\nAblOdomCommon}{-2.5}
\newcommand{\nAblOdomComplex}{+1.7}
\newcommand{\nAblOdomVisual}{-2.9}
\newcommand{\nAblOdomLong}{+1.3}
\newcommand{\nAblOdomOtherMin}{-2.9}
\newcommand{\nAblOdomOtherMax}{+2.1}
\newcommand{\nAblRgbglobal}{-26.4}
\newcommand{\nAblRgbglobalAbs}{26.4}
\newcommand{\nAblRgbglobalLo}{-32.4}
\newcommand{\nAblRgbglobalHi}{-20.5}
\newcommand{\nAblRgbglobalP}{{<}10^{-5}}
\newcommand{\nAblRgbglobalBase}{-32.9}
\newcommand{\nAblRgbglobalCommon}{-25.8}
\newcommand{\nAblRgbglobalComplex}{-28.3}
\newcommand{\nAblRgbglobalVisual}{-16.2}
\newcommand{\nAblRgbglobalLong}{-28.7}
\newcommand{\nAblRgbglobalOtherMin}{-32.9}
\newcommand{\nAblRgbglobalOtherMax}{-16.2}
\newcommand{\nAblRgbfloor}{-2.9}
\newcommand{\nAblRgbfloorAbs}{2.9}
\newcommand{\nAblRgbfloorLo}{-6.3}
\newcommand{\nAblRgbfloorHi}{0.3}
\newcommand{\nAblRgbfloorP}{{=}0.67}
\newcommand{\nAblRgbfloorBase}{-4.6}
\newcommand{\nAblRgbfloorCommon}{-3.3}
\newcommand{\nAblRgbfloorComplex}{-2.5}
\newcommand{\nAblRgbfloorVisual}{-0.4}
\newcommand{\nAblRgbfloorLong}{-3.8}
\newcommand{\nAblRgbfloorOtherMin}{-4.6}
\newcommand{\nAblRgbfloorOtherMax}{-0.4}
\newcommand{\nBudEBTen}{16.0}
\newcommand{\nBudEBFifteen}{44.7}
\newcommand{\nBudSTen}{28.2}
\newcommand{\nBudSFifteen}{68.5}
\newcommand{\nBudRGBTen}{22.5}
\newcommand{\nBudRGBFifteen}{62.8}
\newcommand{\nBudOracleTen}{40.0}
\newcommand{\nBudOracleFifteen}{87.0}


\title{Ask Small VLMs Only What They Can See:\\A Semantic-to-Metric Agent for Embodied Navigation\\[0.3em]
\large Supplementary Material}
\author{Anonymous submission}

\begin{document}
\maketitle

\section{Implementation details}
\label{sec:s-impl}

\noindent\textbf{Parameters.}
\Cref{tab:s-params} lists every parameter of \ours. All were fixed on the development split (\cref{sec:s-protocol})
and are the same for both models and for the RGB-only configuration; RoboTHOR changes only the items of the
explicit-stopping paragraph of Sec.~4.4 of the main paper.

\begin{table}[t]
\centering\small
\caption{Parameters of \ours on EB-Navigation.}
\label{tab:s-params}
\footnotesize\setlength{\tabcolsep}{3pt}
\begin{tabular}{@{}llr@{}}
\toprule
Group & Parameter & Value\\
\midrule
Questions & presence threshold for named targets & 0.5\\
 & crop side (times the box's longer side) & 2.5\\
 & crop and tile resolution & 448\,px\\
 & crop acceptance & 0.5\\
 & weight of a rejected crop & $0.2\,\pi^z_t$\\
Lifting & depth window (fraction of the box) & 0.5\\
 & maximum range & 6\,m\\
 & floor / obstacle height & 0.10 / 0.1--1.9\,m\\
 & occupancy grid & 5\,cm\\
Memory & association radius $r$ & 0.5\,m\\
 & repeated-view discount $\gamma$ & 0.25\\
 & same view: position, heading & 0.2\,m, equal\\
 & negative evidence $\lambda$ & 0.4\\
 & image margin, occlusion tolerance & 20\,px, 0.3\,m\\
 & removal threshold & $w_k\le-1$\\
Search & commit $\tau_c$ / keep $\tau_k$ & 0.8 / 0.4\\
 & approach a weak hypothesis & $w_k\ge0.35$\\
 & exhaustion radius & 0.9\,m\\
 & exploration goal distance & 0.75--3\,m\\
 & tiles, overlap, boxes checked & $3{\times}3$, 20\%, 4\\
Planning & cost of an unobserved edge point & 0.15\\
 & obstacle inflation (fallback) & 0.20 (0.10)\,m\\
 & goal tolerance & 0.55\,m\\
\bottomrule
\end{tabular}
\end{table}

\noindent\textbf{Prompts.}
The target phrase $s$ comes from one text-only query per episode:
\begin{quote}\small\ttfamily
A robot must navigate to one object in a house. Instruction: "$\langle\ell\rangle$"
Reply with JSON \{"name": "...", "appearance\_only": true or false, "looks": "..."\} where
- name: the most likely specific household object, as its common name in 1-2 words (e.g. "toaster", "trash can",
"laptop"; never a vague word like "device" or "object").
- appearance\_only: true if the instruction describes the object only by colour, shape, material or texture,
without saying what it is or what it is used for.
- looks: the appearance words from the instruction (colour, shape, material), empty if none.
\end{quote}
$s$ is the name, or the appearance words when \texttt{appearance\_only} is true. The two per-frame queries are
\begin{quote}\small\ttfamily
Does this image show $\langle s\rangle$? Answer yes or no.\\[0.3em]
Locate the $\langle s\rangle$ in the image and output its bounding box as JSON: [\{"bbox\_2d": [x1, y1, x2, y2],
"label": "$\langle s\rangle$"\}]. If it is not in the image, output [].
\end{quote}
The presence query generates one token and reads $P(\textit{yes})$ and $P(\textit{no})$ from the ten most likely
first tokens (\cref{eq:s-presence}). The crop check asks the presence query and each tile the box query, both on
the magnified image. Box replies are parsed with one regular expression; a box narrower than 2\,px counts as no box.
\begin{equation}
\pi_t=\frac{\sum_{y\in\textit{yes}}P(y)}{\sum_{y\in\textit{yes}}P(y)+\sum_{n\in\textit{no}}P(n)} ,
\label{eq:s-presence}
\end{equation}
where the sums run over the case variants of each answer among the ten tokens; $\pi_t=0.5$ if neither appears.

\noindent\textbf{Serving.}
Both models run in vLLM with their 4-bit AWQ checkpoints, temperature 0, and at most 96 generated tokens for a box
and one for a presence answer. Qwen2.5-VL receives frames resized to at most 409{,}600 pixels; Qwen3-VL uses its
default processor. The two per-frame queries are sent in parallel.

\section{Protocol}
\label{sec:s-protocol}

\noindent\textbf{Development split.}
EmbodiedBench provides target objects and the three instruction variants for 90 iTHOR rooms, and its test set uses
60 of them (FloorPlan11--30, 211--230, 311--330). We rerun the benchmark's own generator on the other 30
(FloorPlan1--10, 201--210, 301--310): a start pose 2.5--3.0\,m from the first instance of the target with the camera
level, Common, Complex, and Visual from the same targets with the benchmark's instruction variants, and Long as Base
with the start heading rotated by 180$^\circ$. This yields five subsets of 30 episodes. Every parameter of
\cref{tab:s-params} and every design choice of Sec.~4 was selected on this split, and the configuration was
frozen before the first test run.

\noindent\textbf{Statistics.}
All agents run on the same 300 test episodes, so each comparison uses per-episode differences. A configuration run
several times enters with its per-episode mean success. Confidence intervals are percentile intervals of $10^4$
episode-level bootstrap resamples. $p$-values are two-sided sign-flip permutation tests with $10^6$ random sign
assignments; the ten variants of Tab.~2 of the main paper are Holm-corrected together. Comparisons of two single
binary runs use the exact McNemar test.

\noindent\textbf{Outcome categories.}
A failed episode takes the first matching category: the target never enters the view; the final distance is at most
1.5\,m (near miss); the target is out of view in the last five steps (lost after seen); at least 30\% of the actions
are rejected (blocked); at least three immediate reversals (oscillation); otherwise approach failure. Visibility
comes from the simulator's object metadata.

\noindent\textbf{RoboTHOR.}
\nRoboErr{} of the \nRoboN{} validation episodes end because the simulator rejects the start pose
(\texttt{teleport failed}); they are counted as failures.

\section{Additional results}
\label{sec:s-results}

\noindent\textbf{Single views.}
\Cref{tab:s-probe} extends the probe of Sec.~3 of the main paper to both models and to the views that do not show
the target. Qwen3-VL-4B finds the target more often than Qwen2.5-VL-3B but also places more boxes where there is
nothing to find, most of all for appearance-only phrases. Both kinds of error enter the memory, and both are what
multi-view fusion and negative evidence act on.

\begin{table*}[t]
\centering\small
\caption{Single-view probe on the test episodes of the first four subsets (start view and the three other 90$^\circ$
headings; 75 views with the target visible per subset).}
\label{tab:s-probe}
\begin{tabular}{@{}lcccccccc@{}}
\toprule
& \multicolumn{4}{c}{Box on the target, IoU${\ge}0.3$ (\%)} & \multicolumn{4}{c}{Box on a view without it (\%)}\\
\cmidrule(lr){2-5}\cmidrule(l){6-9}
Model & Base & Common & Complex & Visual & Base & Common & Complex & Visual\\
\midrule
Qwen3-VL-4B & 78.7 & 69.3 & 77.3 & 49.3 & 14.5 & 18.8 & 16.4 & 41.2\\
Qwen2.5-VL-3B & 40.0 & 33.3 & 37.3 & 18.7 & 3.6 & 8.5 & 3.0 & 8.5\\
\bottomrule
\end{tabular}

\end{table*}

\noindent\textbf{Outcomes.}
\Cref{tab:s-taxonomy} gives the outcome of every agent on the 300 test episodes. \ours removes almost all
never-seen and blocked episodes with both models. With Qwen2.5-VL-3B the EB prompt loses a target it has seen in
\nLostEBTwo{} episodes; \ours does so in \nLostSTwo, its largest remaining failure mode with the smaller model.

\begin{table}[t]
\centering\small
\caption{Outcomes on EB-Navigation test (episodes of 300; mean over runs where a configuration was run
repeatedly).}
\label{tab:s-taxonomy}
\footnotesize\setlength{\tabcolsep}{3pt}
\begin{tabular}{@{}lrrrrr@{}}
\toprule
& \multicolumn{3}{c}{Qwen3-VL-4B} & \multicolumn{2}{c}{Qwen2.5-VL-3B}\\
\cmidrule(lr){2-4}\cmidrule(l){5-6}
Outcome & EB prompt & \ours & \ours, RGB & EB prompt & \ours\\
\midrule
success & 165 & 238.2 & 229.5 & 61 & 210\\
never seen & 12 & 0.8 & 0 & 14 & 0\\
lost after seen & 43 & 28.5 & 24.5 & 104 & 48\\
near miss & 34 & 9.8 & 20.5 & 26 & 20\\
blocked & 26 & 0.5 & 1 & 51 & 3\\
oscillation & 1 & 10.2 & 9 & 1 & 2\\
approach failure & 19 & 12 & 15.5 & 43 & 17\\
\bottomrule
\end{tabular}

\end{table}

\noindent\textbf{Cost.}
Per episode, \ours makes \nCallsS{} VLM calls against \nCallsEB{} for the EB prompt, but each call asks one short
question: it reads \nTokS{}k prompt tokens and generates \nTokSc{} tokens, against \nTokEB{}k and \nTokEBc{}k for the EB
prompt, whose plans and reasoning dominate its generation.

\noindent\textbf{RoboTHOR by category.}
\Cref{tab:s-robothor} lists success and SPL for each of the twelve categories.

\begin{table}[t]
\centering\small
\caption{\ours on RoboTHOR validation by target category.}
\label{tab:s-robothor}
\begin{tabular}{@{}lrrr@{}}
\toprule
Category & Episodes & SR & SPL\\
\midrule
basketball & 150 & 88.0 & 54.3\\
mug & 150 & 77.3 & 45.7\\
apple & 150 & 74.7 & 50.1\\
garbage can & 150 & 71.3 & 47.2\\
television & 150 & 62.7 & 38.5\\
baseball bat & 150 & 62.7 & 34.4\\
laptop & 150 & 61.3 & 27.8\\
house plant & 150 & 56.0 & 28.7\\
spray bottle & 150 & 49.3 & 31.8\\
alarm clock & 150 & 44.7 & 29.5\\
bowl & 150 & 41.3 & 23.8\\
vase & 150 & 25.3 & 18.5\\
\midrule
all & 1800 & 59.6 & 35.9\\
\bottomrule
\end{tabular}

\end{table}

\end{document}
